\documentclass[12pt,letterpaper]{article}
\usepackage{graphicx,textcomp}
\usepackage{natbib}
\usepackage{setspace}
\usepackage[margin=1in]{geometry}
\usepackage{color}
\usepackage[reqno]{amsmath}
\usepackage{amsthm}
\usepackage{fancyvrb}
\usepackage{amssymb,enumerate}
\usepackage[all]{xy}
\usepackage{endnotes}
\usepackage{lscape}
\newtheorem{com}{Comment}
\usepackage{float}
\usepackage{hyperref}
\newtheorem{lem} {Lemma}
\newtheorem{prop}{Proposition}
\newtheorem{thm}{Theorem}
\newtheorem{defn}{Definition}
\newtheorem{cor}{Corollary}
\newtheorem{obs}{Observation}
\usepackage[compact]{titlesec}
\usepackage{dcolumn}
\usepackage{tikz}
\usetikzlibrary{arrows}
\usepackage{multirow}
\usepackage{xcolor}
\newcolumntype{.}{D{.}{.}{-1}}
\newcolumntype{d}[1]{D{.}{.}{#1}}
\definecolor{light-gray}{gray}{0.65}
\usepackage{url}
\usepackage{listings}
\usepackage{color}

\definecolor{codegreen}{rgb}{0,0.6,0}
\definecolor{codegray}{rgb}{0.5,0.5,0.5}
\definecolor{codepurple}{rgb}{0.58,0,0.82}
\definecolor{backcolour}{rgb}{0.95,0.95,0.92}

\lstdefinestyle{mystyle}{
	backgroundcolor=\color{backcolour},   
	commentstyle=\color{codegreen},
	keywordstyle=\color{magenta},
	numberstyle=\tiny\color{codegray},
	stringstyle=\color{codepurple},
	basicstyle=\footnotesize,
	breakatwhitespace=false,         
	breaklines=true,                 
	captionpos=b,                    
	keepspaces=true,                 
	numbers=left,                    
	numbersep=5pt,                  
	showspaces=false,                
	showstringspaces=false,
	showtabs=false,                  
	tabsize=2
}
\lstset{style=mystyle}
\newcommand{\Sref}[1]{Section~\ref{#1}}
\newtheorem{hyp}{Hypothesis}

\title{Problem Set 2}
\date{Due: February 18, 2026}
\author{Applied Stats II}


\begin{document}
	\maketitle
	\section*{Instructions}
	\begin{itemize}
		\item Please show your work! You may lose points by simply writing in the answer. If the problem requires you to execute commands in \texttt{R}, please include the code you used to get your answers. Please also include the \texttt{.R} file that contains your code. If you are not sure if work needs to be shown for a particular problem, please ask.
		\item Your homework should be submitted electronically on GitHub in \texttt{.pdf} form.
		\item This problem set is due before 23:59 on Wednesday February 18, 2026. No late assignments will be accepted.

	\end{itemize}

	\vspace{.25cm}

\noindent We're interested in what types of international environmental agreements or policies people support (\href{https://www.pnas.org/content/110/34/13763}{Bechtel and Scheve 2013)}. So, we asked 8,500 individuals whether they support a given policy, and for each participant, we vary the (1) number of countries that participate in the international agreement and (2) sanctions for not following the agreement. \\

\noindent Load in the data labeled \texttt{climateSupport.RData} on GitHub, which contains an observational study of 8,500 observations.

\begin{itemize}
	\item
	Response variable: 
	\begin{itemize}
		\item \texttt{choice}: 1 if the individual agreed with the policy; 0 if the individual did not support the policy
	\end{itemize}
	\item
	Explanatory variables: 
	\begin{itemize}
		\item
		\texttt{countries}: Number of participating countries [20 of 192; 80 of 192; 160 of 192]
		\item
		\texttt{sanctions}: Sanctions for missing emission reduction targets [None, 5\%, 15\%, and 20\% of the monthly household costs given 2\% GDP growth]
		
	\end{itemize}
	
\end{itemize}

\newpage
\noindent Please answer the following questions:

\begin{enumerate}
	\item
	Remember, we are interested in predicting the likelihood of an individual supporting a policy based on the number of countries participating and the possible sanctions for non-compliance.
	\begin{enumerate}
		\item [] Fit an additive model. Provide the summary output, the global null hypothesis, and $p$-value. Please describe the results and provide a conclusion.

	\end{enumerate}
	
	\item
	If any of the explanatory variables are statistically significant in this model, then:
	\begin{enumerate}
		\item
		For the policy in which nearly all countries participate [160 of 192], how does increasing sanctions from 5\% to 15\% change the odds that an individual will support the policy? (Interpretation of a coefficient)
		\item
		For the policy in which very few countries participate [20 of 192], how does increasing sanctions from 5\% to 15\% change the odds that an individual will support the policy? (Interpretation of a coefficient)
		\item
		What is the estimated probability that an individual will support a policy if there are 80 of 192 countries participating with no sanctions? 
	
	\end{enumerate}
	\item
	Would the answers to 2a and 2b potentially change if we included an interaction term in this model? Why? 
	\begin{itemize}
		\item Perform a test to see if including an interaction is appropriate.
	\end{itemize}
	\end{enumerate}
	\newpage
	\section*{Solutions}
	\vspace{.5cm}
	\begin{enumerate}
		\item 
		\begin{itemize}
			\item \textbf{Fitting the additive model :}
				\lstinputlisting[language=R, firstline=39,lastline=53]{PS02_DP.R} 
			\item \textbf{Summary output: }
			\begin{table}[!htbp] \centering 
				\caption{} 
				\label{} 
				\begin{tabular}{@{\extracolsep{5pt}}lc} 
					\\[-1.8ex]\hline 
					\hline \\[-1.8ex] 
					& \multicolumn{1}{c}{\textit{Dependent variable:}} \\ 
					\cline{2-2} 
					\\[-1.8ex] & choice \\ 
					\hline \\[-1.8ex] 
					countries (80 of 192) & 0.336$^{***}$ \\ 
					& (0.054) \\ 
					& \\ 
					countries (160 of 192) & 0.648$^{***}$ \\ 
					& (0.054) \\ 
					& \\ 
					5\% sanctions & 0.192$^{***}$ \\ 
					& (0.062) \\ 
					& \\ 
					15\% sanctions & $-$0.133$^{**}$ \\ 
					& (0.062) \\ 
					& \\ 
					20\% sanctions & $-$0.304$^{***}$ \\ 
					& (0.062) \\ 
					& \\ 
					Constant & $-$0.273$^{***}$ \\ 
					& (0.054) \\ 
					& \\ 
					\hline \\[-1.8ex] 
					\textit{Note:}  & \multicolumn{1}{r}{$^{*}$p$<$0.1; $^{**}$p$<$0.05; $^{***}$p$<$0.01} \\ 
				\end{tabular} 
			\end{table} 
			\item \textbf{Global null hypothesis test:} \\[0.25cm]
				The global null hypothesis test compares the null deviance (deviance of the null model) with the residual model (of the full additive model), using a Likelihood Ratio Test (LRT). The null hypothesis is that all slope coefficients are equal to 0, and the alternative hypothesis is that at least one of the slopes is different from 0.
				\lstinputlisting[language=R, firstline=55,lastline=61]{PS02_DP.R} 
				\begin{verbatim}
					Analysis of Deviance Table
					
					Model 1: choice ~ 1
					Model 2: choice ~ countries + sanctions
					Resid. Df Resid. Dev Df Deviance  Pr(>Chi)    
					1      8499      11783                          
					2      8494      11568  5   215.15 < 2.2e-16 ***
					---
					Signif. codes:  0 ‘***’ 0.001 ‘**’ 0.01 ‘*’ 0.05 ‘.’ 0.1 ‘ ’ 1
				\end{verbatim}
			\item \textbf {Interpretation and Conclusion:} \\[0.25cm]
				Looking at the output of the test, we can see that the p-value is much lower than 0.05. Thus, we can reject the null hypothesis that all coefficients are equal to 0. This shows that at least one of the predictors in the full model is a reliable one in the logistic regression. The full model does a significantly better job at predicting whether an individual supports the policy or not (the outcome variable), than the null model does. 
		\end{itemize}
		\newpage
		\item \textbf{Refitting the model - changing the reference category for the sanctions variable from None to 5\% : }
		\lstinputlisting[language=R, firstline=67,lastline=75]{PS02_DP.R} 
		\begin{table}[!htbp] \centering 
			\caption{} 
			\label{} 
			\begin{tabular}{@{\extracolsep{5pt}}lc} 
				\\[-1.8ex]\hline 
				\hline \\[-1.8ex] 
				& \multicolumn{1}{c}{\textit{Dependent variable:}} \\ 
				\cline{2-2} 
				\\[-1.8ex] & choice \\ 
				\hline \\[-1.8ex] 
				countries (80 of 192) & 0.336$^{***}$ \\ 
				& (0.054) \\ 
				& \\ 
				countries (160 of 192) & 0.648$^{***}$ \\ 
				& (0.054) \\ 
				& \\ 
				no sanctions & $-$0.192$^{***}$ \\ 
				& (0.062) \\ 
				& \\ 
				15\% sanctions & $-$0.325$^{***}$ \\ 
				& (0.062) \\ 
				& \\ 
				20\% sanctions & $-$0.495$^{***}$ \\ 
				& (0.062) \\ 
				& \\ 
				Constant & $-$0.081 \\ 
				& (0.053) \\ 
				& \\ 
				\hline \\[-1.8ex] 
				\textit{Note:}  & \multicolumn{1}{r}{$^{*}$p$<$0.1; $^{**}$p$<$0.05; $^{***}$p$<$0.01} \\ 
			\end{tabular} 
		\end{table} 
			\lstinputlisting[language=R, firstline=77,lastline=77]{PS02_DP.R} 
			\begin{verbatim}
				0.7224531 
			\end{verbatim}
			Looking at the stars next to the coefficients, we can infer that all the explanatory variables are statistically significant in this model. Thus, we can proceed with the interpretation of different coefficients, as the associations they capture are unlikely to be due to chance. 
		\begin{enumerate}
			\item \textbf{Interpretation:} \\[0.25cm]
			 Increasing sanctions from 5\% to 15\%, decreases the log-odds that an individual will support the policy by 0.325, holding the number of participating countries constant. Thus, within a given category of participating countries, the odds of an individual supporting a policy with 15\% sanctions are 0.72 times the odds of supporting a policy with 5\% sanctions. There is a 28\% decrease in the odds of someone supporting the policy when sanctions increase from 5\% to 15\%, for the policy in which nearly all countries participate.
			 \vspace{0.5cm}
			\item \textbf{Interpretation:} \\[0.25cm]
			Because our model has no interaction term (it is an additive model), an increase from 5\% to 15\% in sanctions has the same effect regardless of the number of participating countries, so it is the same for the policy group in which nearly all countries participate as it is in the group in which very few countries participate. Therefore, the interpretation above (2a) remains unchanged for the policy group in which very few countries participate.  There is also a 28\% decrease in the odds of someone supporting a policy when the sanctions increase from 5\% to 15\%. 
			\vspace{0.5cm}
			\item \textbf{Calculation: }
			\lstinputlisting[language=R, firstline=79,lastline=84]{PS02_DP.R} 
			\begin{verbatim}
				0.5159191 
			\end{verbatim}
		The estimated probability that an individual will support a policy if there are 80 of 192 countries participating with no sanctions is 51.59\%. 
		\end{enumerate}
		\newpage
		\item 
		Yes, the answers to 2a and 2b would potentially change if we included an interaction term in the model. An interaction term would allow us to account for different effects of the sanctions at different levels of countries participating. In the additive model fitted above, the change in log-odds when moving from 5\% to 15\% sanctions is assumed to be the same regardless of how many countries support the policy. Nevertheless, including an interaction term would allow the slope of sanctions to vary depending on the number of participating countries. Increasing sanctions from 5\% to 15\% could actually result in different log-odds of an individual supporting the policy based on whether nearly all countries are participating or very few are. 
		\vspace{0.25cm}
		\begin{itemize}
			\item \textbf{Include an interaction term in the model: }
				\lstinputlisting[language=R, firstline=90,lastline=96]{PS02_DP.R} 
			*** only kept the interaction terms' coefficients to look at their significance
			\begin{table}[!htbp] \centering 
				\caption{} 
				\label{} 
				\begin{tabular}{@{\extracolsep{5pt}}lc} 
					\\[-1.8ex]\hline 
					\hline \\[-1.8ex] 
					& \multicolumn{1}{c}{\textit{Dependent variable:}} \\ 
					\cline{2-2} 
					\\[-1.8ex] & choice \\ 
					\hline \\[-1.8ex] 
					countries (80 of 192) : no sanctions & $-$0.095 \\ 
					& (0.152) \\ 
					& \\ 
					countries (160 of 192) : no sanctions & $-$0.130 \\ 
					& (0.151) \\ 
					& \\ 
					countries (80 of 192) : 15\% sanctions & $-$0.147 \\ 
					& (0.154) \\ 
					& \\ 
					countries (160 of 192) : 15\% sanctions & $-$0.182 \\ 
					& (0.151) \\ 
					& \\ 
					countries (80 of 192) : 20\% sanctions & $-$0.292$^{*}$ \\ 
					& (0.153) \\ 
					& \\ 
					countries (160 of 192) : 20\% sanctions & $-$0.073 \\ 
					& (0.152) \\ 
					& \\ 
					\hline \\[-1.8ex] 
					\textit{Note:}  & \multicolumn{1}{r}{$^{*}$p$<$0.1; $^{**}$p$<$0.05; $^{***}$p$<$0.01} \\ 
				\end{tabular} 
			\end{table} 
			\newpage
			\item \textbf{Perform test: }
			\vspace{.25cm}
				\lstinputlisting[language=R, firstline=98,lastline=99]{PS02_DP.R} 
				\begin{verbatim}
					Analysis of Deviance Table
					
					Model 1: choice ~ countries + sanctions
					Model 2: choice ~ countries * sanctions
					Resid. Df Resid. Dev Df Deviance Pr(>Chi)
					1      8494      11568                     
					2      8488      11562  6   6.2928   0.3912
				\end{verbatim}
			\item \textbf{Conclusion: }\\[0.25cm]
			The p-value is higher than the 0.05 threshold, which shows that there is not enough evidence that including an interaction term is a significant predictor for the odds of an individual supporting the policy. The additive model does a good enough job at predicting these odds. This conclusion could have also been reached by looking at the individual regression coefficients, as we can see that none of the interaction coefficients is statistically reliable. This also suggests that the effect of increasing sanctions on the outcome does not depend on the number of participating countries, so the change in log-odds associated with sanctions is similar across different levels of country participation.
		\end{itemize}
	\end{enumerate}

\end{document}
